% !TeX root = ../../../main.tex
\subsubsection{Survival Analysis Models}
\label{ch03:sec:methods_survival}
% \vspace{-0.25em}
To estimate the adverse event time, we consider four survival analysis models that vary in their underlying assumptions and modeling flexibility. 
These models include both traditional statistical approaches and more recently developed deep learning methods and can be plugged into meta-learning causal inference approaches:
% Below, we describe the four survival models used in our experiments:
\begin{enumerate}\setlength{\itemsep}{0pt}
    \item \textbf{Cox Proportional Hazards (CoxPH)} \citep{cox1972regression} – A semiparametric model that assumes a log-linear relationship between covariates and the hazard function, with a time-invariant hazard ratio:
    \(
    h(t \mid \mathbf{Z}_{i\tau}) = h_0(t) \exp(\boldsymbol{\beta}^{\top} \mathbf{Z}_{i\tau}),
    \)
    where \( h_0(t) \) is the baseline hazard and \( \boldsymbol{\beta} \) the regression coefficients.
CoxPH is widely used for its interpretability and computational efficiency.

    \item \textbf{Random Survival Forest (RSF)} \citep{ishwaran2008random} – A nonparametric ensemble method that extends decision trees to survival data.
It constructs multiple survival trees and aggregates them to estimate the survival function:
    \(
    \widehat{S}(t \mid \mathbf{Z}_{i\tau}) = \frac{1}{B} \sum_{b=1}^{B} S^{(b)}(t \mid \mathbf{Z}_{i\tau}),
    \)
    where \( S^{(b)}(t \mid \mathbf{Z}_{i\tau}) \) is the survival function from the \( b \)-th tree.
RSF captures non-linear relationships and complex interactions without parametric assumptions.

    \item \textbf{DeepSurv} \citep{katzman2018deepsurv} – A deep learning extension of CoxPH that replaces its linear assumption with a neural network:
    \(
    h(t \mid \mathbf{Z}_{i\tau}) = h_0(t) \exp(f_{\theta}(\mathbf{Z}_{i\tau})),
    \)
    where \( f_{\theta} \) is a neural network.
DeepSurv learns non-linear covariate effects while maintaining the proportional hazards framework.

    \item \textbf{DeepHit} \citep{lee2018deephit} – A neural network-based model that directly estimates the probability mass function (PMF) of survival time:
    \(
    P(T = t \mid \mathbf{Z}_{i\tau}) = f_{\theta}(t, \mathbf{Z}_{i\tau}),
    \)
    where \( f_{\theta} \) predicts the likelihood of event occurrence at each time \( t \).
DeepHit allows for modeling multimodal survival distributions and competing risks.
\end{enumerate}
Our experiments later also uses a Causal Survival Forest (CSF) \citep{cui2023estimating}, which explicitly accounts for treatment effects in survival analysis.
However, since a CSF is inherently a causal model that estimates treatment effects, it is not meant for being plugged into a meta-learning causal inference framework.
We defer discussing it to the next section, where we describe the causal inference setup we consider. 

% \vspace{-1em}
